\begin{center}
    \textbf{\Huge{Relatório de Spike Técnico - Estudo de Caracterização e Integração de Modelos de Inteligência Artificial}}
\end{center}

\section{Introdução}

O propósito deste documento é apresentar os resultados de um spike técnico realizado para aprofundar o entendimento sobre os modelos de Inteligência Artificial atualmente utilizados no projeto, bem como sobre a forma como esses modelos são estruturados, executados e como podem ser eventualmente integrados ao simulador. O trabalho foi motivado pela necessidade de reduzir incertezas técnicas e estabelecer uma base conceitual mais sólida para decisões de engenharia relacionadas à arquitetura da solução.

\subsection{Referências}

\begin{itemize}
    \item Documento n\textsuperscript{o} ASA2/MIIA/001, ``\textbf{Visão de Produto}'', Instituto de Estudos Avançados (IEAv), 2025
    \item Repositório ASA2/MIIA, Merge Request 10, \url{https://gitlab.asa.dcta.mil.br/asa/asa2-miia/-/commit/9811d09abd5180fac9fac4a377f6987685d9e009}, Instituto de Estudos Avançados (IEAv), 2026
\end{itemize}
\newpage

\section{Caracterização dos Modelos de IA}

Esta seção apresenta a caracterização técnica dos modelos de Inteligência Artificial utilizados no projeto, com foco em sua estrutura, funcionamento e implicações para a arquitetura do sistema. O objetivo é descrever de forma concreta o que esses modelos representam do ponto de vista de engenharia de software e como se dá sua interação com o ambiente de simulação.

\subsection{Estrutura e Representação dos Modelos}

Os modelos de IA analisados são produtos de treinamento realizado no ecossistema da biblioteca \emph{Stable-Baselines3}, em particular utilizando o algoritmo \emph{PPO} (\emph{Proximal Policy Optimization}). Do ponto de vista prático, esses modelos não correspondem a um único arquivo executável ou artefato isolado, mas a um conjunto de arquivos e estruturas que encapsulam os pesos da rede neural, metadados do treinamento, configuração do ambiente e dependências do \emph{framework} utilizado. É comum que essa estrutura esteja organizada na forma de arquivos \emph{ZIP}, com nomes indicando o tipo de algoritmo utilizado, \emph{seed} e número de \emph{steps} de treinamento para uma versão específica do modelo.

A seguir apresenta-se a estrutura típica de um modelo treinado com PPO no formato gerado pela biblioteca Stable-Baselines3:

\vspace{0.5cm}

\begin{lstlisting}[language=bash]
  $ less PPO_seed_1975_3600000_steps.zip
  Archive:  PPO_seed_1975_3600000_steps.zip
 Length   Method    Size  Cmpr    Date    Time   CRC-32   Name
--------  ------  ------- ---- ---------- ----- --------  ----
   21353  Stored    21353   0% 2025-12-11 15:23 6b8183bc  data
    1261  Stored     1261   0% 1980-01-01 00:00 6cdd3008  pytorch_variables.pth
  156156  Stored   156156   0% 1980-01-01 00:00 048b9650  policy.pth
  313006  Stored   313006   0% 1980-01-01 00:00 56116f1c  policy.optimizer.pth
       5  Stored        5   0% 2025-12-11 15:23 51d40e63  _stable_baselines3_version
     241  Stored      241   0% 2025-12-11 15:23 4063029b  system_info.txt
--------          -------  ---                            -------
  492022           492022   0%                            6 files
\end{lstlisting}

\vspace{0.5cm}

A execução direta desses modelos depende do ambiente \emph{Python} e das bibliotecas que os suportam, o que impõe restrições relevantes em termos de portabilidade, versionamento e integração com outros componentes do sistema que não compartilham o mesmo ecossistema de execução. Por exemplo, a forma de carregamento de um modelo treinado com o algoritmo \emph{PPO} é diferente de um modelo treinado com o algoritmo \emph{DQN} (\emph{Deep Q-Network}), conforme ilustrado a seguir:

\vspace{0.5cm}

\begin{lstlisting}[language=python]
  from stable_baselines3 import PPO, DQN

  ...
  
  # API Proximal Policy Optimization
  model = PPO.load(model_path, device="cpu")
  # API Deep Q-Network
  model = DQN.load(model_path, device="cpu")
\end{lstlisting}

O que prejudica a generalidade de eventuais soluções de integração desses modelos, haja vista que não apenas restringiria o uso da linguagem de programação para Python mas também demandaria o uso de condicionais específicos a depender de como o modelo foi treinado.

\subsection{Conceitos Fundamentais: State, Observations e Actions}

Durante a análise, tornou-se necessário estabelecer uma distinção precisa entre os conceitos centrais que regem a interação entre o simulador e os modelos de Inteligência Artificial: \emph{state}, \emph{observation} e \emph{action}.

O \textit{state} corresponde à representação completa do cenário simulado em um dado instante. Ele é mantido internamente pelo simulador e contém todas as informações relevantes do ambiente, como posições das entidades, velocidades, orientações, estados internos de sensores, entre muitos outros atributos físicos e lógicos.

A \textit{observation}, por sua vez, não é uma cópia direta do \textit{state}. Ela consiste em uma visão parcial e transformada desse estado, construída especificamente para atender às necessidades de um determinado modelo de IA. Em geral, a \textit{observation} utiliza apenas um subconjunto dos atributos do \textit{state} e aplica etapas de pré-processamento, como normalização e cálculo de grandezas relativas.

A normalização consiste na transformação de valores físicos de diferentes escalas (por exemplo, metros, graus, metros por segundo) para uma faixa numérica comum, tipicamente entre 0 e 1 ou entre -1 e 1, com o objetivo de facilitar o processo de aprendizado do modelo.

Um exemplo de componentes de uma \textit{observation} é apresentado a seguir, ilustrando como o modelo recebe não valores brutos do simulador, mas grandezas relativas e normalizadas para duas entidades simuladas B1 e R1:

\begin{table}[h!]
    \centering
    \begin{tabular}{|l|}
        \hline
        \textbf{\emph{Observation}} \\
        \hline
        Diferença de magnitude normalizada dos azimutes relativos para (B1,R1) \\
        \hline
        Distância oblíqua normalizada para (B1,R1) \\
        \hline
        Diferença de magnitude normalizada das velocidades aerodinâmicas para (B1,R1) \\
        \hline
        Diferença de magnitude normalizada dos azimutes relativos para (B1,R2) \\
        \hline
        Distância oblíqua normalizada para (B1,R2) \\
        \hline
        Diferença de magnitude normalizada das velocidades aerodinâmicas para (B1,R2) \\
        \hline
        Diferença de magnitude normalizada dos azimutes relativos para (B2,R1) \\
        \hline
        Distância oblíqua normalizada para (B2,R1) \\
        \hline
        Diferença de magnitude normalizada das velocidades aerodinâmicas para (B2,R1) \\
        \hline
        Diferença de magnitude normalizada dos azimutes relativos para (B2,R2) \\
        \hline
        Distância oblíqua normalizada para (B2,R2) \\
        \hline
        Diferença de magnitude normalizada das velocidades aerodinâmicas para (B2,R2) \\
        \hline
        Diferença de magnitude normalizada dos azimutes relativos para (B1,B2) \\
        \hline
        Distância oblíqua normalizada para (B1,B2) \\
        \hline
        Diferença de magnitude normalizada das velocidades aerodinâmicas para (B1,B2) \\
        \hline
    \end{tabular}
\end{table}

O termo \textit{action} corresponde à saída produzida pelo modelo a partir da \textit{observation} recebida. Essa ação representa a decisão tomada pela política aprendida durante o treinamento e é utilizada pelo simulador para atualizar o \textit{state}, influenciando diretamente a evolução do cenário.

A relação entre \textit{state}, \textit{observation} e \textit{action} define o ciclo fundamental de interação entre o simulador e o modelo de IA.


\subsection{Ciclo de Inferência e Integração com o Sistema}

O ciclo de inferência pode ser descrito, de forma simplificada, pelas seguintes etapas:
\begin{itemize}
    \item O simulador mantém e atualiza continuamente o \textit{state} completo do ambiente.
    \item A partir do \textit{state}, constrói-se a \textit{observation} específica para o modelo de IA, aplicando seleção de atributos e pré-processamento (por exemplo, normalização e cálculo de grandezas relativas).
    \item A \textit{observation} é fornecida ao modelo de IA como vetor de entrada.
    \item O modelo executa a inferência de forma determinística e produz uma \emph{action}.
    \item A \textit{action} é aplicada ao ambiente pelo sistema, modificando o \textit{state}.
    \item O ciclo se repete no próximo passo de simulação.
\end{itemize}

Esse fluxo estabelece uma dependência clara entre a representação interna do simulador e a interface esperada pelo modelo, criando um acoplamento técnico que deve ser cuidadosamente gerenciado do ponto de vista arquitetural.


\section{Desenvolvimento do Conversor SB3 PPO \textrightarrow{} ONNX}

\subsection{Motivação}

Durante a caracterização dos modelos de IA utilizados no projeto, tornou-se evidente que a dependência direta do ecossistema Python e das bibliotecas específicas de treinamento impunha limitações significativas à arquitetura do sistema. A execução de inferência estava fortemente acoplada às mesmas ferramentas utilizadas durante o treinamento, dificultando a portabilidade dos modelos, o versionamento controlado e a integração com componentes externos ao ambiente Python.

Além disso, a ausência de um formato padronizado e independente para a representação dos modelos impedia a utilização de runtimes especializados de inferência e restringia as possibilidades de implantação do sistema em diferentes ambientes computacionais. Diante desse cenário, tornou-se necessário investigar uma estratégia que permitisse desacoplar o modelo treinado do seu ambiente original de execução, preservando seu comportamento funcional e ampliando sua interoperabilidade.

Nesse contexto, a conversão dos modelos PPO treinados com a biblioteca Stable-Baselines3 para o formato \emph{ONNX} foi adotada como abordagem inicial, como estudo de viabilidade para uma arquitetura mais modular, portátil e sustentável.

\subsection{Implementação}

\subsubsection{Estratégia adotada}

A estratégia consistiu em isolar explicitamente a política treinada pelo algoritmo \emph{PPO} e exportá-la para o formato \emph{ONNX}, de modo que o modelo resultante pudesse ser executado por \emph{runtimes} de inferência independentes do ambiente de treinamento original. Esse processo envolveu o encapsulamento da política em um módulo \emph{PyTorch}, o carregamento dos pesos treinados e a exportação da rede para o formato \emph{ONNX} por meio das ferramentas disponibilizadas pelo próprio \emph{framework}.

A inferência foi configurada de forma determinística, assegurando reprodutibilidade e permitindo a validação numérica direta entre os resultados produzidos no ambiente original e aqueles obtidos pelo modelo exportado.

\subsubsection{Ferramentas utilizadas}

Foram utilizadas as seguintes ferramentas e tecnologias no desenvolvimento do conversor:

\begin{itemize}
    \item Stable-Baselines3 para carregamento do modelo PPO treinado;
    \item PyTorch para reconstrução e manipulação da rede neural;
    \item Utilitários de exportação ONNX fornecidos pelo PyTorch;
    \item ONNX Runtime para validação e execução de inferência do modelo exportado.
\end{itemize}

\subsubsection{Desafios técnicos}

Durante a implementação, foram enfrentados desafios relacionados à correta reconstrução da política de decisão, à compatibilidade entre operações suportadas pelo ONNX e aquelas utilizadas no modelo original, bem como à validação da equivalência funcional entre a inferência executada no ambiente original e a inferência executada via runtime externo.

Esses desafios exigiram ajustes na definição da rede, tratamento cuidadoso de tensores de entrada e saída e validações sistemáticas do comportamento do modelo convertido.

\subsection{Resultados}

O processo de conversão resultou em um modelo funcional no formato \emph{ONNX}, capaz de executar inferência de forma independente do ambiente de treinamento original em \emph{Python}. O modelo exportado pôde ser executado com sucesso utilizando um \emph{runtime} externo, mantendo equivalência numérica com o modelo original treinado em \emph{Stable-Baselines3}, conforme verificado por testes sistemáticos de erro máximo e médio.

Essa abordagem proporcionou benefícios diretos à arquitetura do sistema, incluindo maior portabilidade dos modelos, redução do acoplamento com bibliotecas de treinamento, facilitação do versionamento e ampliação das possibilidades de integração com outros componentes e plataformas de execução.

\begin{lstlisting}[language=bash]
  $ sb3ppo2onnx PPO_seed_1975_3600000_steps.zip PPO_seed_1975_3600000_steps.onnx
  Modelo carregado: PPO_seed_1975_3600000_steps.zip
  Modelo exportado para ONNX: PPO_seed_1975_3600000_steps.onnx
  
  Verificando equivalencia PyTorch <---> ONNX
  Teste 1: max error = 0.000000 | mean error = 0.000000
  Teste 2: max error = 0.000000 | mean error = 0.000000
  Teste 3: max error = 0.000000 | mean error = 0.000000
  Teste 4: max error = 0.000000 | mean error = 0.000000
  Teste 5: max error = 0.000000 | mean error = 0.000000
  Verificacao concluida: exportacao valida.
\end{lstlisting}

Cabe ressaltar que a adoção do formato \emph{ONNX}, no contexto deste trabalho, deve ser entendida como um estudo técnico de viabilidade e não como uma decisão arquitetural definitiva. A avaliação realizada teve como objetivo verificar a possibilidade de desacoplamento entre o ambiente de treinamento e o ambiente de execução, bem como investigar os impactos dessa estratégia sobre portabilidade, desempenho e integridade funcional do modelo. A decisão final quanto ao formato de representação e ao \emph{runtime} de inferência a ser adotado pelo sistema será tomada em etapas posteriores do projeto, à luz dos resultados obtidos neste estudo e de outros requisitos técnicos e operacionais ainda em análise.
